\documentclass[11pt,letterpaper]{article}
% Shared declarations; root documents always import a concrete semantic module.
\usepackage{business-admin}
\usepackage{amsmath,amssymb}
\usepackage{longtable,tabularx,array,booktabs}
\usepackage{enumitem}
\usepackage{ragged2e}
\newcolumntype{Y}{>{\RaggedRight\arraybackslash}X}
\newcommand{\SuiteMeta}[2]{\def\DocTitle{#1}\def\DocSubtitle{#2}\def\DocVersion{1.0}\def\DocStatus{Educational reference}\def\DocOwner{Small Business Management Suite}\def\DocSummary{Original teaching and practice resource; no real-business approval implied.}}
\newcommand{\SuiteOpen}{\makedoctitle\vspace{-0.7em}\noindent\rule{\linewidth}{0.45pt}\par\medskip}
\newcommand{\SourceNote}[2]{\par\medskip\noindent{\footnotesize\textbf{Reading anchor:} Longenecker, Petty, Palich and Hoy, \emph{Small Business Management: Launching \& Growing Entrepreneurial Ventures}, 20th ed., Ch.~#1, textbook p.~#2. This document is independently authored commentary and exercises, not a reproduction of the chapter.}}
\newcommand{\WritingLines}[1]{\par\smallskip\noindent\rule{\linewidth}{0.35pt}\par\vspace{0.35cm}\noindent\rule{\linewidth}{0.35pt}\par\vspace{0.35cm}\noindent\rule{\linewidth}{0.35pt}\par\medskip}
\newcommand{\FormField}[1]{\par\medskip\noindent\textbf{#1}\par\WritingLines{3}}
\newcommand{\SBMCheckBox}{\(\square\)}
\newcommand{\AssessmentRow}[2]{#1 & #2 & \(\square\) & \(\square\) & \(\square\) & \(\square\) & \(\square\)\\\hline}
\newcommand{\AssessmentHeader}{\noindent\begin{tabular}{|p{0.22\linewidth}|p{0.49\linewidth}|c|c|c|c|c|}\hline\textbf{Criterion} & \textbf{Evidence to inspect} & \textbf{0} & \textbf{1} & \textbf{2} & \textbf{3} & \textbf{4}\\\hline}
\newcommand{\AssessmentEnd}{\end{tabular}\par}
\setlength{\emergencystretch}{3em}

\SuiteMeta{Chapter 17: Promotional Planning}{Study guide | entrepreneurship and small business management}
\begin{document}
\SuiteOpen
\section{Learning outcome}
Design a promotion program whose spending is tied to observable customer outcomes.
\section{Prerequisites}
Read chapters 7, 14, 16.\par\textbf{Required baseline:} ability to record assumptions and distinguish observations from opinions.
\section{Chapter concept map}
\begin{itemize}[leftmargin=*,itemsep=3pt]
\item promotion objectives
\item budget methods
\item website and social channels
\item personal selling
\item advertising and sales promotion
\end{itemize}

\section{Technical study notes}
\begin{itemize}[leftmargin=*,itemsep=3pt]
\item Promotion is a communication system with a target audience, message, channel, budget and observable behavior, not a list of advertising tactics.
\item Objective-and-task budgeting connects spend to an outcome but still requires measurement and constraints. Percent-of-sales budgeting may ignore opportunity quality.
\item Reach, engagement, leads and purchases reflect different funnel stages; avoid evaluating success with a metric disconnected from contribution or retention.
\end{itemize}

\section{Reading and analysis protocol}
\begin{enumerate}[leftmargin=*,itemsep=5pt]
\item Read the chapter learning objectives, then mark terms and claims that need primary evidence.
\item For each chapter theme, distinguish the entrepreneur decision, information required, parties affected and failure condition.
\item Write a one-paragraph synthesis in your own words; cite the textbook chapter and separate contemporary practice notes.
\item Complete the applied activity and attach the underlying evidence, calculations or interview log.
\end{enumerate}

\section{Worked application and deliverable}
\textbf{Business exercise.} Create a channel experiment card with a capped budget and stop rule.\par\textbf{Required submission.} Promotion experiments and budget. Include a dated assumption register and an explicit proceed/revise/stop recommendation.\par\textbf{Assurance checks.} Attribution, offer accuracy, acquisition cost and measurement.
\section{Common reasoning failures}
\begin{itemize}[leftmargin=*,itemsep=3pt]
\item Optimizing impressions without revenue
\item Ignoring seasonality or attribution limits
\end{itemize}
\noindent\textbf{Correction:} require counterevidence and identify the responsible decision-maker before approving the next stage.
\section{Discussion prompt}
Why might a campaign with positive click-through still be rejected?\par\medskip\textbf{Evidence expected:} at least two alternatives, one quantitative or independently checkable observation, a trade-off, and a decision threshold.
\section{Exit assessment}
\begin{enumerate}[leftmargin=*,itemsep=5pt]
\item Explain each term from the concept map without copying the textbook.
\item Defend the applied decision using an assumption, evidence and counterargument.
\item Show an objective test or inspection that could prove a key assumption wrong.
\item State what would change the recommendation and who has approval authority.
\end{enumerate}

\SourceNote{17}{441}
\end{document}
